% Imporditud: tools/import_eksperimendid.py
% Allikas: Eesti füüsikaolümpiaadi eksperimendi ülesannete
%   kogu (Rael Kalda, 2023), ülesanne Ü102, lahendus L102.
\setAuthor{EFO žürii}
\setRound{lõppvoor}
\setYear{2011}
\setNumber{P E2}
\setDifficulty{3}
\setTopic{termodünaamika}
\setVahendid{kalorimeeter, termomeeter, mõõtejoonlaud, lumi, vesi}
\setPunktid{12}
\setAllikas{Eksperimendi kogu Ü102, lahendus lk 79}
\setLahendusseis{toores}

\prob{Lumi}
Määrake lume sulamissoojuse ja vee erisoojuse suhe.

\hint

\solu
Täidame kalorimeetri umbes kolmveerandi ulatuses veega ning mõõdame joonlauaga veetaseme kõrguse $h_1$ ja vee algtemperatuuri $T_1$. Nüüd paneme vette järjest lund ja segame seda joonlauaga lume täieliku sulamiseni, kuni veetase on sulanud lume tõttu märgatavalt tõusnud ning veetemperatuur langenud. Mõõdame taas vee temperatuuri $T_2$ ja kõrguse $h_2$. Leiame nüüd seosed avaldamaks lume sulamissoojuse $\lambda$ ja vee erisoojuse c suhte. Algselt on kalorimeetris oleva vee mass $m_1$ = Sh1$\rho$, kus S on kalorimeetri ristlõikepindala ja $\rho$ vee tihedus. Lume massi m$^2$ saame leida ruumalamuudust: m$^2$ = ($h_2$ $-$ $h_1$) S$\rho$. Eeldades, et välise keskkonnaga soojusvahetus puudub, kuna kalorimeeter on isoleeritud, saame kirjutada soojushulkade kohta võrrandi:

m1c ($T_1$ $-$ $T_2$) = $\lambda$m$^2$ + m2c ($T_2$ $-$ 0) .

Asendades võrrandisse $m_1$ ja m$^2$ ning lihtsustades saame otsitud suuruse:

$\lambda$ = $h_1$ ($T_1$ $-$ $T_2$) $-$ $T_2$ = h1T1 $-$ h2T2 . c $h_2$ $-$ $h_1$ $h_2$ $-$ $h_1$

Vastuseks on $\lambda$ $\approx$ 63 $^\circ$C. c
\probend
